\documentclass[11pt]{article}

\usepackage[T1]{fontenc}
\usepackage[utf8]{inputenc}
\usepackage{lmodern}
\usepackage{amsmath,amssymb,mathtools}
\usepackage{booktabs}
\usepackage{dsfont}
\usepackage[hidelinks]{hyperref}

\title{TmFeO$_3$ Hamiltonian and Tm Magnetization Readout}
\author{}
\date{\today}

\DeclareMathOperator{\diag}{diag}
\newcommand{\bS}{\mathbf{S}}
\newcommand{\bh}{\mathbf{h}}
\newcommand{\be}{\mathbf{e}}
\newcommand{\bD}{\mathbf{D}}
\newcommand{\Jeff}{\mathbf{J}^{\rm eff}}
\newcommand{\deff}{\mathbf{d}^{\rm eff}}
\newcommand{\half}{\tfrac{1}{2}}

\begin{document}

\maketitle

This note summarizes the Hamiltonian implemented in
\texttt{src/core/unitcell\_builders.cpp} for \texttt{build\_tmfeo3()},
using the notation and physical definitions of the companion theory
document \texttt{tmfeo3\_notes.tex}. Section~\ref{sec:readout} records
how the effective Tm SU(3) variables are transformed when magnetization
is saved, and Section~\ref{sec:drive} details the time-dependent drive.

\tableofcontents

% -------------------------------------------------------------------
\section{Conventions}
\label{sec:conventions}

\subsection{Sites and sublattice frames}

The mixed TmFeO$_3$ unit cell contains four Fe sites and four Tm sites.
Fe spins are three-component SU(2) vectors $\bS_i = (S_i^x, S_i^y,
S_i^z)$ with spin length $S$; Tm sites are stored as eight-component
real Gell-Mann expectation-value vectors
$\boldsymbol{\lambda}_j = (\lambda_1, \dots, \lambda_8)$ in the local
crystal-field (CEF) basis.

The four Pbnm sublattices are related by the Klein four-group of proper
$\pi$-rotations $\{R_\mu\}_{\mu=0}^{3} = \{E, C_{2x}, C_{2y}, C_{2z}\}$,
represented as diagonal matrices
\begin{equation}
R_0 = \diag(+,+,+),\quad R_1 = \diag(+,-,-),\quad
R_2 = \diag(-,+,-),\quad R_3 = \diag(-,-,+).
\end{equation}
In the code, the same matrices are written as
$D_\mu = \diag(\eta_{\mu,x}, \eta_{\mu,y}, \eta_{\mu,z})$.  A global
spin is recovered from the local-frame spin by
$S^\alpha_\mu\big|_{\rm glob} = (R_\mu)_{\alpha\alpha}\,\tilde S^\alpha_\mu\big|_{\rm loc}$.

\subsection{Bertaut modes}

The four sublattice combinations
\begin{equation}
\mathbf{X}_{\mathbf{n}} = \tfrac14\sum_{\mu=0}^{3} \sigma^X_\mu\,\bS_{\mathbf{n},\mu},
\quad X \in \{F,G,C,A\},
\end{equation}
with sign vectors
\begin{equation}
\sigma^F=(+,+,+,+),\;\sigma^G=(+,-,-,+),\;
\sigma^C=(+,+,-,-),\;\sigma^A=(+,-,+,-),
\end{equation}
span the Bertaut basis.  The inverse relation is
$\bS_{\mathbf{n},\mu} = \sum_X \sigma^X_\mu\,\mathbf{X}_\mathbf{n}$.
The $\Gamma_2$ ground state has $\mathbf{G} = G_z\hat z$,
$\mathbf{F} = F_x\hat x$, $\mathbf{C} = C_y\hat y$, $\mathbf{A}\approx 0$.

\subsection{Gell-Mann classification}

The eight Gell-Mann generators are classified by mirror parity ($m_z$,
site symmetry $C_s$) and time-reversal parity $T$:
\begin{equation}
\boxed{
\begin{aligned}
A_1^+\ (\text{mirror-even,}\ T\text{-even}):&\quad\lambda^1,\lambda^3,\lambda^8,\\
A_1^-\ (\text{mirror-even,}\ T\text{-odd}):&\quad\lambda^2,\\
A_2^+\ (\text{mirror-odd,}\ T\text{-even}):&\quad\lambda^4,\lambda^6,\\
A_2^-\ (\text{mirror-odd,}\ T\text{-odd}):&\quad\lambda^5,\lambda^7.
\end{aligned}}
\label{eq:lambdaclass}
\end{equation}
The $T$-odd generators $\lambda^2, \lambda^5, \lambda^7$ carry
the projected magnetic dipole $\Jeff$; the $T$-even off-diagonal
generators $\lambda^1, \lambda^4, \lambda^6$ carry the projected electric
dipole $\deff$.

\subsection{Projected dipole operators}

In the $T$-symmetric gauge the three lowest Tm CEF singlets
$\{|E_1(A_1)\rangle, |E_2(A_1)\rangle, |E_3(A_2)\rangle\}$ have
numerical matrix elements (see \texttt{tmfeo3\_notes.tex} \S III.C)
\begin{equation}
\Jeff = g_J\mu_B\begin{pmatrix}
J_z^{\rm eff} \\ J_x^{\rm eff} \\ J_y^{\rm eff}
\end{pmatrix}
= g_J\mu_B\begin{pmatrix}
+5.264\,\lambda^2 \\
+2.392\,\lambda^5 + 0.913\,\lambda^7 \\
-2.787\,\lambda^5 + 0.466\,\lambda^7
\end{pmatrix},
\label{eq:Jeff}
\end{equation}
which defines the $3\times 8$ active-channel matrix $\mu_{\rm act}$
used in the code:
\begin{equation}
\mu_{\rm act} =
\begin{pmatrix}
0 & \mu_{2x} & 0 & 0 & \mu_{5x} & 0 & \mu_{7x} & 0 \\
0 & \mu_{2y} & 0 & 0 & \mu_{5y} & 0 & \mu_{7y} & 0 \\
0 & \mu_{2z} & 0 & 0 & \mu_{5z} & 0 & \mu_{7z} & 0
\end{pmatrix},
\end{equation}
so that $(J^\alpha)^{\rm eff} = \sum_{a\in\{2,5,7\}}\mu_{\alpha a}\lambda^a$.
The projected electric dipole $\deff$ populates $\lambda^1$ (via $d_{x,y}$)
and $\lambda^{4,6}$ (via $d_z$).

% -------------------------------------------------------------------
\section{Fe Sector}
\label{sec:Fe}

\subsection{Four-sublattice Hamiltonian}

The Fe sector is the standard four-sublattice orthoferrite Hamiltonian
with Heisenberg exchange, Dzyaloshinskii--Moriya (DM) interaction, and
biaxial single-ion anisotropy:
\begin{equation}
\begin{aligned}
H_{\rm Fe} = {}&
\sum_{\langle ij\rangle_{\parallel}} J_1\,\bS_i\!\cdot\!\bS_j
 + \sum_{\langle ij\rangle_c} J_{1c}\,\bS_i\!\cdot\!\bS_j
 + \sum_{\langle\!\langle ij\rangle\!\rangle} J_2\,\bS_i\!\cdot\!\bS_j\\
 &+\sum_{\langle ij\rangle_{\parallel}} \bD_{ij}\cdot(\bS_i\times\bS_j)\\
 &- \sum_i \bigl[K_a(S_i^x)^2 + K_b(S_i^y)^2 + K_c(S_i^z)^2\bigr].
\end{aligned}
\label{eq:HFe}
\end{equation}
Here $J_1$ (\texttt{J1ab}), $J_{1c}$ (\texttt{J1c}), $J_2$
(\texttt{J2ab}/\texttt{J2c}) are the in-plane NN, c-axis NN, and NNN
exchange constants; $K_{a,b,c}$ (\texttt{Ka}, \texttt{Kb}, \texttt{Kc})
are the biaxial single-ion anisotropy constants.

\subsection{DM bond pattern}

The DM vectors on in-plane NN bonds are fixed by $Pbnm$ symmetry.
With bonds at height $z = \tfrac{1}{2}$ and $z = 0$ distinguished by
the $S_1$ screw axis:
\begin{equation}
\bD^{(z=\tfrac{1}{2})}_{ij} = (0,+D_1,+D_2),\qquad
\bD^{(z=0)}_{ij} = (0,-D_1,+D_2).
\label{eq:Dpattern}
\end{equation}
The sign reversal of $D_1$ between the two planes is required by
the $S_1$ screw that maps one $z$-plane to the other; $D_2$ is
symmetric.

\subsubsection*{Code implementation}

The code encodes each bond as a lab-frame exchange matrix via
\begin{equation}
J(J_{\rm iso}, D_y, D_z) =
\begin{pmatrix}
J_{\rm iso} & D_z & -D_y \\
-D_z & J_{\rm iso} & 0 \\
D_y & 0 & J_{\rm iso}
\end{pmatrix}.
\end{equation}
With the parameter names in \texttt{apply\_tmfeo3\_fe\_sector()}, the
NN exchange matrices assigned to each bond set are
\begin{align}
J_a^{(10)} = J_b^{(10)} &= J(J_{1ab}, +D_1, +D_2), \notag\\
J_a^{(23)} = J_b^{(23)} &= J(J_{1ab}, -D_1, +D_2), \\
J_c,\; J_{2a},\; J_{2b},\; J_{2c} &= J_{\rm iso}\,\mathds{1}_{3\times 3}.\notag
\end{align}
The antisymmetric part of $J_{ij}$ is identified with
$\bD_{ij}\cdot(\bS_i\times\bS_j)$, reproducing Eq.~\eqref{eq:Dpattern}.

\subsection{Fe Zeeman (static field)}

A static external field $\bh = h\,\hat n$ adds
\begin{equation}
H_{\rm Fe,Z} = -g_e\mu_B\sum_i \bh\cdot\bS_i,
\end{equation}
where in the code the $g_e\mu_B$ prefactor is absorbed into
\texttt{field\_strength} $= g_e\mu_B\,h$ and
\texttt{field\_direction} $= \hat n$.

% -------------------------------------------------------------------
\section{Tm Sector}
\label{sec:Tm}

\subsection{CEF Hamiltonian: three-level truncation}

The microscopic Tm$^{3+}$ CEF Hamiltonian is a Stevens-operator
sum over even-$|m|$ terms allowed by the $C_s$ site symmetry,
truncated to the lowest three singlets
$\{|E_1(A_1)\rangle, |E_2(A_1)\rangle, |E_3(A_2)\rangle\}$ with
energies $\epsilon_1 = 0$, $\epsilon_2 \equiv e_1$,
$\epsilon_3 \equiv e_2$ (code parameters \texttt{e1}, \texttt{e2}).
In the truncated qutrit basis the diagonal part is
\begin{equation}
H_{\rm CEF}^{(0)} = \bar\epsilon\,\mathds{1}
+ h_3^{(0)}\lambda^3 + h_8^{(0)}\lambda^8,
\label{eq:HTm0}
\end{equation}
with
\begin{equation}
h_3^{(0)} = \frac{e_1}{2},\qquad
h_8^{(0)} = \frac{e_1 - 2e_2}{2\sqrt{3}},
\end{equation}
and $\bar\epsilon = (e_1 + e_2)/3$ irrelevant.  Numerically,
$h_3^{(0)} = -0.96$\,meV, $h_8^{(0)} = -3.97$\,meV.  The code stores
the on-site SU(3) field vector
\begin{equation}
\mathbf{B}_{\rm CEF} = (0, 0, \alpha, 0, 0, 0, 0, \beta),
\end{equation}
with
\begin{equation}
\alpha = -h_3^{(0)}\,s_\alpha,\qquad
\beta  = -h_8^{(0)}\,s_\beta,
\end{equation}
and scale factors \texttt{tm\_alpha\_scale} $= s_\alpha$,
\texttt{tm\_beta\_scale} $= s_\beta$ (default $1$).
The on-site Tm CEF contribution per site is
$H_{\rm CEF}^{(0)} = -\mathbf{B}_{\rm CEF}\cdot\boldsymbol\lambda_j$.

\subsection{Tm Zeeman coupling (static field)}

The static field couples to the projected magnetic dipole:
\begin{equation}
H_{\rm Tm,Z} = -\sum_j \Jeff_j \cdot \bh
= -g_{\rm ratio}\sum_j\sum_{a\in\{2,5,7\}} B_a\,\lambda_a^{(j)},
\end{equation}
where $B_a = \sum_\alpha \mu_{\alpha a}\,h_\alpha$ projects the field
onto the active $T$-odd channels, and
\texttt{g\_ratio\_tm} $= g_{\rm ratio} = g_J/g_e$ is the
Tm-to-Fe $g$-factor ratio.

\subsection{Tm--Tm diagonal bilinears}

Nearest-neighbor Tm--Tm interactions are parameterized as a diagonal
coupling in the Gell-Mann basis:
\begin{equation}
H_{\rm TmTm} = \frac{1}{2}\sum_{\langle j,j'\rangle}\sum_{a=1}^{8}
J_a^{\rm Tm}\,\lambda_a^{(j)}\lambda_a^{(j')},
\end{equation}
implemented via the $8\times 8$ diagonal matrix
$J_{\rm Tm} = \diag(J_{\rm tm,1}, \dots, J_{\rm tm,8})$ on the
bond geometry listed in \texttt{apply\_tmfeo3\_tm\_sector()}.

% -------------------------------------------------------------------
\section{Fe--Tm Couplings}
\label{sec:FeTm}

\subsection{Bilinear Fe--Tm coupling ($\chi$ term)}

Because the Hamiltonian must be $T$-even and the Fe spin $\bS$ is
$T$-odd, only the $T$-odd Gell-Mann generators $\lambda^{2,5,7}$
can appear linearly (Eq.~\eqref{eq:lambdaclass}).  The most general
$Pbnm$-symmetric bilinear Fe--Tm coupling is
\begin{equation}
H_\chi = \sum_{(i,j)}\sum_{\alpha=x,y,z}\sum_{a\in\{2,5,7\}}
\chi^{\alpha a}_{ij}\,S_i^\alpha\,\lambda_j^a,
\label{eq:Hchi}
\end{equation}
with bond-resolved coupling tensor $\chi^{\alpha a}_{ij}$.  The
active channels in matrix form are
\begin{equation}
\chi =
\begin{pmatrix}
0 & \chi_{2x} & 0 & 0 & \chi_{5x} & 0 & \chi_{7x} & 0 \\
0 & \chi_{2y} & 0 & 0 & \chi_{5y} & 0 & \chi_{7y} & 0 \\
0 & \chi_{2z} & 0 & 0 & \chi_{5z} & 0 & \chi_{7z} & 0
\end{pmatrix}.
\end{equation}
At $\mathbf{q}=0$ only $\bar\chi^{2z}\neq 0$ survives the
inversion-pair cancellation; this is the dominant channel coupling
$\lambda^2$ ($A_1^-$) to $G_z$ (the $\Gamma_2$ N\'{e}el order
parameter).  The off-diagonal mirror-odd channels
$\chi^{5x,7x,5y,7y}$ cancel at $\mathbf{q}=0$ due to the Pbnm
bond-pair structure.

The code stores 16 explicit inversion-related bond pairs;
on the inversion-odd partner bond the $\lambda^5$ and $\lambda^7$
channels flip sign while $\lambda^2$ does not.

\subsection{Mixed trilinear $W$ term}

The $T$-even generators $\lambda^{1,3,4,6,8}$ do not appear in the
linear coupling $H_\chi$, but they enter at second order via virtual
Tm transitions.  Expanding to quadratic order in the Bertaut
amplitudes, the resulting $\mathbf{q}=0$ effective Hamiltonian is
\begin{equation}
\begin{aligned}
H_W = {}&
{-}\sum_{\mathbf{n},j}\bigl[
u_1\lambda^1_{\mathbf{n},j}
+ u_3\lambda^3_{\mathbf{n},j}
+ u_8\lambda^8_{\mathbf{n},j}\bigr]
\mathcal{I}^{\rm Fe}_{A_1^+}(\mathbf{n})\\
&-\sum_{\mathbf{n},j}\bigl[
v_4\lambda^4_{\mathbf{n},j}
+ v_6\lambda^6_{\mathbf{n},j}\bigr]
\mathcal{I}^{\rm Fe}_{A_2^+}(\mathbf{n}),
\end{aligned}
\label{eq:HW}
\end{equation}
where the $A_1^+$ (mirror-even) Fe bilinears are
\begin{equation}
\mathcal{I}^{\rm Fe}_{A_1^+}
= \alpha_1 G_z^2 + \alpha_2 F_x^2 + \alpha_3 C_y^2 + \cdots
\end{equation}
and the $A_2^+$ (mirror-odd) ones are
\begin{equation}
\mathcal{I}^{\rm Fe}_{A_2^+}
= \beta_1 F_xG_z + \beta_2 G_zC_y + \cdots
\end{equation}
with $\alpha_i, \beta_i$ phenomenological $\mathbf{q}=0$ coefficients.
In the code, \texttt{build\_tmfeo3()} populates
$n \in \{1,3,4,6,8\}$ channels of the on-site trilinear tensor
$W^{\alpha\beta}_{n,ij}$.
The $A_2^+$ channels ($\lambda^4$, $\lambda^6$) change sign on
inversion-odd bond partners; the $A_1^+$ channels ($\lambda^{1,3,8}$)
are even.  The $A_2^+$ on-site contribution vanishes at $\mathbf{q}=0$
in the inversion-pair cancellation, leaving the $A_1^+$
population-modulation channel as the dominant back-action mechanism.

\subsection{Total Hamiltonian}

The full working Hamiltonian is
\begin{equation}
\boxed{
H_{\rm total}(t) = H_{\rm Fe} + H_{\rm CEF}^{(0)}
+ H_{\rm Tm,Z} + H_{\rm TmTm}
+ H_\chi + H_W
+ H_{\rm drive}(t),
}
\label{eq:Htotal}
\end{equation}
corresponding to the microscopic Hamiltonian $H_{\rm micro}(t)$
of \texttt{tmfeo3\_notes.tex} Eq.~(Hmicro) augmented by the $W$ back-action
term, Eq.~(Htotal).

% -------------------------------------------------------------------
\section{How the Tm Magnetization Is Saved}
\label{sec:readout}

\subsection{Stored local SU(3) variables}

Each Tm site is evolved and stored internally as an eight-component
local vector $\boldsymbol{\lambda}^{\rm loc} = (\lambda_1, \dots,
\lambda_8)$ in the canonical local CEF basis.

\subsection{Sublattice-frame map for readout}

The builder constructs an $8\times 8$ sublattice frame $F_\mu$ from
the projected dipole matrix:
\begin{equation}
F_\mu = \mu_{\rm act}^{-1}\,D_\mu\,\mu_{\rm act}
\quad\text{on the active triplet }\{\lambda_2,\lambda_5,\lambda_7\}.
\end{equation}
The same $3\times 3$ action is extended to the coherence-partner
triplet $\{\lambda_1,\lambda_4,\lambda_6\}$, while $\lambda_3$ and
$\lambda_8$ are left invariant.  In block form,
\begin{equation}
\boldsymbol{\lambda}^{\rm glob} = F_\mu\,\boldsymbol{\lambda}^{\rm loc}.
\end{equation}
The frame is installed via
\texttt{Tm\_atoms.set\_sublattice\_frame(\ldots)} and copied into
\texttt{sublattice\_frames\_SU3} when the \texttt{MixedLattice} is built.

\subsection{What the saved SU(3) magnetization means}

When the code computes global SU(3) magnetization it applies the frame
site by site:
\begin{equation}
M_{{\rm SU3},\mu}^{\rm glob} =
\frac{1}{N_{\rm Tm}}\sum_{j\in\mu} F_\mu\,\boldsymbol{\lambda}_j^{\rm loc}.
\end{equation}
\begin{itemize}
\item \texttt{M\_local\_SU3}: average local qutrit Bloch vector in the
  CEF basis.
\item \texttt{M\_global\_SU3}: average SU(3) vector after $F_\mu$ is
  applied.
\item \texttt{M\_antiferro\_SU3}: staggered combination formed by the
  mixed-lattice observable code.
\end{itemize}

\subsection{Recovering the physical Tm dipole}

The physically observable Tm magnetic dipole in lab Cartesian
components is
\begin{equation}
J^\alpha_{\rm lab}(j) = \eta_{\mu(j),\alpha}\sum_{a\in\{2,5,7\}}
\mu_{\alpha a}\,\lambda_a^{\rm loc}(j).
\end{equation}
The SU(3) arrays are not themselves the Cartesian dipole; they are
the Gell-Mann expectation values in the local CEF basis or after the
sublattice-frame map.  The physical Tm magnetization is recovered by
projecting the $T$-odd channels with $\mu_{\rm act}$ and then applying
the Pbnm sign pattern $D_\mu$.

% -------------------------------------------------------------------
\section{Drive Hamiltonian}
\label{sec:drive}

The THz drive couples to both Fe and Tm via magnetic-dipole and
(at Tm) electric-dipole vertices.  The microscopic drive terms are
(\texttt{tmfeo3\_notes.tex} Eqs.~(Hdrivem)--(Hdrivee))
\begin{align}
H_{\rm drive}^{\rm m}(t) &= -g_e\mu_B\sum_{i}\bh(t)\cdot\bS_i
- g_J\mu_B\sum_{j}\bh(t)\cdot\Jeff_j, \label{eq:Hdm}\\
H_{\rm drive}^{\rm e}(t) &= -\sum_{j}\be(t)\cdot\deff_j. \label{eq:Hde}
\end{align}
The two THz propagation geometries are (for propagation along $\hat b$):
\begin{equation}
\begin{aligned}
H\!\parallel\! a,\; E\!\parallel\! c:\;&
h_x\neq 0,\; e_z\neq 0
\;\Rightarrow\;\text{drives }\lambda^{5,7}\text{ and }\lambda^{4,6};\\
H\!\parallel\! c,\; E\!\parallel\! a:\;&
h_z\neq 0,\; e_x\neq 0
\;\Rightarrow\;\text{drives }\lambda^2\text{ and }\lambda^1.
\end{aligned}
\end{equation}

\subsection{Pulse envelope}

The code supports up to two simultaneously active pulses ($k = 1, 2$)
per species.  Each pulse carries a Gaussian-modulated cosine envelope:
\begin{equation}
f_k(t) =
A\exp\!\left[-\frac{(t-t_{B_k})^2}{4\sigma^2}\right]
\cos\!\bigl[\omega(t-t_{B_k})\bigr],
\end{equation}
where $A$ (\texttt{pump\_amplitude}), $\sigma$ (\texttt{pump\_width}),
and $\omega$ (\texttt{pump\_frequency}) are set independently for the
Fe (SU(2)) and Tm (SU(3)) sectors.  The intensity profile
$I(t)\propto e^{-(t-t_{B_k})^2/(2\sigma^2)}$ has standard deviation
$\sigma$.  A tabulated-waveform mode (via \texttt{pump\_table\_file})
replaces the Gaussian with linear interpolation of a normalized
$E(t)$ file.

\subsection{Fe drive}

The pulse direction for Fe atom type $\mu$ is a 3-component Cartesian
vector $\mathbf{b}^{(k)}_\mu$, which is transformed to the local
sublattice frame before integration:
\begin{equation}
\tilde{\mathbf{b}}^{(k)}_\mu = D_\mu^T\,\mathbf{b}^{(k)}_\mu.
\end{equation}
The drive contribution to the local effective field on Fe site $i$ is
\begin{equation}
\mathbf{h}^{\rm Fe}_{{\rm drive},i}(t)
= \sum_{k=1}^{2} f_k(t)\,\tilde{\mathbf{b}}^{(k)}_\mu,
\end{equation}
giving the Hamiltonian term
$H_{\rm drive}^{\rm m,Fe}(t) = -\sum_i \mathbf{h}^{\rm Fe}_{{\rm drive},i}(t)\cdot\bS_i$,
reproducing the $-g_e\mu_B\sum_i\bh\cdot\bS_i$ term of
Eq.~\eqref{eq:Hdm} with $A = g_e\mu_B\,h_{\rm peak}$.

\subsection{Tm drive}

The pulse direction for Tm atom type $\nu$ is an 8-component Gell-Mann
vector $\mathbf{c}^{(k)}_\nu$, transformed to the local CEF frame as
\begin{equation}
\tilde{\mathbf{c}}^{(k)}_\nu = F_\nu^T\,\mathbf{c}^{(k)}_\nu.
\end{equation}
The drive contribution to the local SU(3) effective field on Tm site $j$
is
\begin{equation}
\mathbf{h}^{\rm Tm}_{{\rm drive},j}(t)
= \sum_{k=1}^{2} f_k(t)\,\tilde{\mathbf{c}}^{(k)}_\nu,
\end{equation}
giving
$H_{\rm drive}^{\rm Tm}(t) = -\sum_j
\mathbf{h}^{\rm Tm}_{{\rm drive},j}(t)\cdot\boldsymbol{\lambda}_j$.

\paragraph{Channel assignments for physical fields.}
Any of the eight Gell-Mann channels may be driven by setting the
corresponding component of \texttt{pump\_directions\_su3}.  For a
physical electromagnetic wave with field components $h_\alpha(t)$ and
$e_\alpha(t)$:
\begin{itemize}
\item \emph{Magnetic} $h_z$ $\Rightarrow$ $c_2 = g_{\rm ratio}\mu_{2z}$
  \quad[$\lambda^2$, $A_1^-$, drives $\omega_{12}$ coherence]
\item \emph{Magnetic} $h_x$ $\Rightarrow$ $c_5 = g_{\rm ratio}\mu_{5x}$,
  $c_7 = g_{\rm ratio}\mu_{7x}$
  \quad[$\lambda^{5,7}$, $A_2^-$, drives $\omega_{13}$/$\omega_{23}$]
\item \emph{Electric} $e_z$ $\Rightarrow$ $c_4 \propto p_{z4}$,
  $c_6 \propto p_{z6}$
  \quad[$\lambda^{4,6}$, $A_2^+$, drives $\omega_{13}$/$\omega_{23}$]
\item \emph{Electric} $e_x$ $\Rightarrow$ $c_1 \propto p_x$
  \quad[$\lambda^1$, $A_1^+$, drives $\omega_{12}$ coherence]
\end{itemize}

\paragraph{The \texttt{auto\_su3\_pump} shortcut.}
When \texttt{auto\_su3\_pump} is enabled, the runner projects a
3-component physical direction $\hat{\mathbf{e}}$ onto the active
$T$-odd channels via $c_a = \sum_\alpha \mu_{\alpha a}\,e_\alpha$,
$a\in\{2,5,7\}$.  Internally this writes to the 0-based array positions
$\{1,4,6\}$, which correspond to $\lambda_2, \lambda_5, \lambda_7$ in
the 1-based Gell-Mann labeling used throughout this document.

\end{document}
